\chapter{The Reversibility Construction}
\label{sec:reversibility}
% ===================================================================

In a bare GSLT, rewrite rules are directed: once a term $P$ reduces to $Q$,
there is in general no canonical way to recover $P$ from $Q$.  This
\emph{irreversibility} is the computational analogue of thermodynamic
time-asymmetry, and a physics needs it undone --- we want a theory whose laws
are time-symmetric and whose arrow of time lives somewhere else.

The construction that does this is already in the book.  It is
Chapter~\ref{ch:histmonad}'s history monad, and this chapter is what that monad
looks like when it is read as physics rather than as bookkeeping.

\section{The envelope is the free history}

\begin{construction}[Reversible envelope]
\label{con:revenv}
Let $\GSLT = (\terms, \eqs, \rules)$.  The \emph{reversible envelope}
$\GSLT^\dagger$ has:
\begin{itemize}[itemsep=2pt]
  \item terms the configurations of $\Hist\GSLT$ --- pairs
        $\langle P, \trace \rangle$ of a current term and a word in the free
        monoid of rewrite events, written here as
        $\trace = r_1(l_1) \cdot r_2(l_2) \cdots r_n(l_n)$ with $l_i$ the
        specific left-hand side matched at step $i$;
  \item forward rules $r^+$, the rewrites of $\Hist\GSLT$:
        $\langle l_r, \trace \rangle \rewrite
         \langle r_r,\, r(l_r) \cdot \trace \rangle$;
  \item backward rules $r^-$, their inverses:
        $\langle r_r,\, r(l_r) \cdot \trace \rangle \rewrite
         \langle l_r, \trace \rangle$;
  \item $\eqs^\dagger$ lifting $\equiv$ to the current-term component, leaving
        traces invariant.
\end{itemize}
\end{construction}

The first two clauses are $\Hist\GSLT$ verbatim.  Only the third is new, and
the point of stating the construction this way is that the third clause is
free.

\begin{proposition}[The backward rules are determined]
\label{prop:bwd-free}
Given the forward rules of $\Hist\GSLT$, the family $\{r^-\}$ is the unique
family of rules inverting them.  Adding it therefore involves no choice.
\end{proposition}

\begin{proof}
By Proposition~\ref{ob:prop:univcover} the free history is the universal cover
of the reduction graph, hence a tree: every configuration
$\langle P, \trace \rangle$ with $\trace \neq \eps$ has exactly one parent, read
off the head of $\trace$.  A rule inverting $r^+$ must send that configuration
to that parent, and $r^-$ does.
\end{proof}

\begin{remark}[What the envelope is, categorically]
\label{rmk:free-groupoid}
Proposition~\ref{prop:bwd-free} says $\GSLT^\dagger$ is the free groupoid on
the tree $\Hist\GSLT$: formally inverting the arrows of a tree imposes no
relations, because there are no loops for the inverses to interact with.
This is the whole reason the reversibility construction is cheap.
Reversibility is not something bought by adding structure to a theory.  It is
what a theory already looks like once it stops throwing its past away, and the
history monad is the operation of stopping.
\end{remark}

Two morphisms come with the monad rather than needing separate proof.  The unit
$\eta : \GSLT \to \GSLT^\dagger$, $P \mapsto \langle P, \eps \rangle$, is
$\Hist$'s unit; the projection $\pi : \GSLT^\dagger \to \GSLT$,
$\langle P, \trace \rangle \mapsto P$, is the functor stripping the log from
Proposition~\ref{ob:prop:hist-adjunction}.  That
$\pi \circ \eta = \mathrm{id}_{\GSLT}$ is one of that adjunction's triangle
identities.

\section{Reading it as physics}

\begin{remark}[Physical Interpretation]
A term $\langle P, \eps \rangle$ is an \emph{initial condition}: it has no causal
past.  A term $\langle P, \trace \rangle$ with $\trace \neq \eps$ is a term with a
recorded history.  The time-asymmetry present in $\GSLT$ (forward rules only)
disappears in $\GSLT^\dagger$: the laws are symmetric, and the arrow of time is
entirely encoded in the boundary condition ($\trace = \eps$ for the initial
state).  This is precisely the modern understanding of the thermodynamic arrow
of time.
\end{remark}

That remark is what this chapter adds, and it is worth being exact about which
part of it the mathematics supplies.  That the laws are symmetric is
Construction~\ref{con:revenv}.  That the asymmetry has to go \emph{somewhere} is
Proposition~\ref{ob:prop:univcover}: a tree has a root, and the root is the
boundary condition.  What is not supplied is any reason the actual world's
boundary condition should be $\trace = \eps$ rather than something else, and
nothing here should be read as offering one.

\begin{remark}[Coarse-graining is an intermediate cover]
\label{rmk:coarse-cover}
$\GSLT^\dagger$ sits at the top of a lattice, and the rungs below it are the
physics.
Chapter~\ref{ch:histmonad} observes that erasures of history are monad
quotients forming a lattice between the reversible cover and the irreversible
base, each intermediate object a quotient of the reduction graph's fundamental
group by a chosen set of loops.
Read physically, an intermediate cover is a \emph{coarse-graining}: a
description that keeps enough of the past to run backwards through some
transitions and not others.
Which loops one is permitted to close is which histories one can afford to
forget --- and by the ledger law $\sigma + \kappa = \sigma_0$ of that chapter,
together with Landauer's principle, the affordable erasures are exactly the
state-recoverable ones.
The fully reversible envelope is the limiting case in which nothing is
forgotten and nothing is paid for.
An observer at an intermediate cover has a time-asymmetric physics not because
the laws are asymmetric but because it cannot afford the receipts.
\end{remark}

\begin{remark}[What is decorated later]
\label{rmk:rev-undecorated}
Nothing in this chapter assigns a value to a rewrite.  Weights, costs and the
semiring in which they live arrive in Chapter~\ref{sec:weighted}, and the
composite object --- terms $\langle P, \trace, \vect{A}\rangle$ carrying a
history \emph{and} an account --- is Construction~\ref{con:main}.
That composite is $\Cost$ applied to $\Hist$, and Chapter~\ref{ch:histmonad}
notes that the order of application is meaningful.  No distributive law between
the two monads is stated anywhere in this book, so ``a cost-accounted theory
with recorded histories'' should be read as the two constructions applied in the
order Construction~\ref{con:main} applies them, and not as a canonical object.
\end{remark}

% ===================================================================
